\SourceLeaf{417}%
Bogens $aa'$, und bezogen auf diesen Punkt bezeichne $\lambda,\Lambda$ dasselbe wie $\alpha,A$ und $\alpha',A'$ in Bezug auf $a$ und $a'$. Ist nun $\alpha\lambda$ oder $\Lambda-A$ dem Bogenstück $al$ proportional, so sagen wir, dass $aa'$ auf seiner ganzen Länge gleichförmig gekrümmt ist, und nennen
\[
\frac{\Lambda-A}{al}
\]
das Krümmungsmaß oder kurz die Krümmung. Man sieht leicht, dass dies nur dann eintritt, wenn $aa'$ wirklich ein Kreisbogen ist, und dass seine Krümmung dann nach unserer Definition $=\pm\dfrac1r$ ist, wenn $r$ den Radius bezeichnet. Fasst man $r$ stets als positiv auf, so gilt das obere oder untere Vorzeichen, je nachdem, ob der Mittelpunkt rechts oder links vom Bogen $aa'$ liegt ($a$ als Anfangspunkt und $a'$ als Endpunkt betrachtet und die Richtungen im Hilfskreis von links nach rechts gezählt). Ändert man eine dieser Bedingungen, so ändert sich das Vorzeichen. Ändert man zwei, so wird das ursprüngliche Vorzeichen wiederhergestellt.

Ist dagegen $\Lambda-A$ nicht proportional zu $al$, so nennen wir den Bogen ungleichförmig gekrümmt. Der Quotient
\[
\frac{\Lambda-A}{al}
\]
kann dann seine mittlere Krümmung heißen. Die Krümmung im eigentlichen Sinn setzt dagegen stets voraus, dass die Stelle bestimmt ist. Sie wird durch die mittlere Krümmung eines Elements an dieser Stelle definiert und ist also
\[
=\frac{d\Lambda}{dal}.
\]
Man sieht also, dass Bogen, Amplitudo und Krümmung in demselben Verhältnis zueinander stehen wie Zeit, Bewegung und Geschwindigkeit oder Volumen, Masse und Dichte. Der Kehrwert der Krümmung, nämlich
\[
\frac{dal}{d\Lambda}
\]
heißt der Krümmungsradius an der Stelle $l$. Unter den obigen Voraussetzungen heißt die Kurve an dieser Stelle nach rechts konkav und nach links konvex, wenn der Wert der Krümmung oder des Krümmungsradius positiv ist. Bei einem negativen Wert verhält es sich umgekehrt.
